\section{Introduction}
\label{sec:intro}

\IEEEPARstart{L}{arge} language model agents increasingly keep a persistent memory. While
triaging tickets, reading documents, or debugging code, an agent records what it has learned and
later retrieves relevant records by nearest-neighbor search over text embeddings
\cite{packer2023memgpt}. The records are written continuously and without human review, and the
memory is often stored on the user's own device. Over time it becomes a detailed account of the
work the agent has performed.

Security research on agent memory has so far concentrated on integrity. Poisoning attacks such
as MINJA \cite{dong2025minja}, AgentPoison \cite{chen2024agentpoison}, and MemoryGraft
\cite{srivastava2025memorygraft} insert records that alter the agent's later behavior, and the
OWASP LLM Top 10 lists vector and embedding weaknesses as a distinct risk
\cite{owasp2025llm08}. Confidentiality has received less attention. MEXTRA \cite{mextra2025}
extracts memory records through the agent's own prompting interface, and membership inference
against chat-agent memory likewise queries the running agent \cite{mrmmia2026}, but no existing
work models an adversary that obtains the storage layer. Encrypting the memory is the natural first
measure, because embeddings can be inverted to recover much of the original text
\cite{morris2023text}.

This paper studies an adversary that obtains the storage layer partway through the life of such a
store, for example through a compromised synchronization or backup account, or by imaging the
device more than once. From the moment of compromise, $\tau$, it observes storage activity and can
influence the agent's inputs; for entries written before $\tau$ it holds only stored state. Those
entries might be expected to be better protected. For a common class of index layouts they are not,
and the reason also says when they are.

\paragraph{The attack, and its precondition}
The attack rests on two properties of the setting. An agent's writes can be induced: an adversary
that supplies a document, a ticket, or an email causes the agent to write a record about it, and
that record is ordinary agent output. And encrypted vector indexes partition their records by
content, so repeated induced writes on a chosen topic reveal which partitions receive records on
that topic. Because partitions are assigned by content rather than by time, the resulting
partition-to-topic map describes the earlier records too. A transfer lemma makes this exact, and
shows the two cohorts are exposed equally.

The lemma has a third hypothesis, and it carries most of the weight: the adversary must learn
\emph{which partition holds a given earlier record}. A layout whose regions grow with occupancy
discloses that in the snapshot; one of fixed-size sealed regions does not. The paper is largely an
account of how an adversary obtains that membership anyway.

\paragraph{What actually protects, and what does not}
Two beliefs about this setting do not survive measurement. The first is that subject drift protects
early records: it lowers the attack's advantage only against an adversary that induces with current
content, and one that induces with archival material recovers it entirely. The second is that a
history-independent representation makes several snapshots no worse than one. History independence
relates a single representation to the operations that produced it and says nothing about a
sequence of representations; two retained versions diff to exactly the regions written between
them, which is the write channel applied before $\tau$.

Two further channels reach the earlier cohort without touching the layout. A store that bounds
per-partition capacity and refuses an overflowing write discloses, at near-perfect accuracy, that a
topic's load has reached capacity, and therefore what the earlier cohort contributed to it. And an
agent that acts on what it recalls leaks through its own operations; the weakest form of this, the
occurrence of a write, is enough, because a memory pipeline that skips a fact it already holds turns
deduplication into a membership oracle for facts written before the compromise.

\paragraph{A design space, not a construction}
The threat model already assumes the adversary cannot read the unlocked process's memory, and in
the local setting there is no remote server that must hold the index. Encrypted on-disk search is
therefore a choice rather than a requirement, and taking the assumption at face value yields a
design whose disk layout carries no content at all: placement is uniform over free slots, search
runs entirely in locked memory, and flushes happen on a clock rather than on writes. That design
closes all three channels above. Its limit is resident memory, and a partition index kept in
locked memory rather than expressed in the disk layout gives it sublinear search at no cost in
leakage. The partitioned design remains the right tool above the memory limit, with its residual
channels reported, including one that is easy to misattribute: its keyed partition permutation is
largely recovered from ordinary search traffic by an adversary that observes reads.

The contributions are as follows.
\begin{enumerate}
\item A threat model for the confidentiality of self-writing agent memory: the $\tau$-arrival
adversary, split into a write-observing and a read-write-observing class, anchored on
synchronization, backup and imaging instantiations, parameterized by version-capture granularity,
and audited against the artifacts a real operating system retains
(\cref{sec:threat,sec:eval-artifacts}).
\item The structure-transfer attack and a transfer lemma that states the condition under which it
reaches the earlier cohort, evaluated with realistic estimators, against deployed stores under
file-level encryption with time-placed negative controls, and against records a language-model
agent actually wrote (\cref{sec:transfer,sec:eval}).
\item Three channels that supply that condition: multi-version reconstruction, the refusal channel
with a capacity--confidentiality proposition, and derived writes, together with a deduplication
membership-inference attack that instantiates the last of them
(\cref{sec:attack-refusal,sec:attack-derived,sec:attack-multisnapshot}).
\item An adaptive, leakage-parameterized definition of dual-epoch security with leakage stated over
plaintext operations, sufficient conditions, a predicate-hiding corollary, and separations realized
as measured ablations (\cref{sec:dualepoch,sec:eval-game}).
\item Three implemented designs spanning the space, with a byte-level leakage auditor that checks
every operation against the bytes it writes and that found a cross-operation defect no functional
test could (\cref{sec:mnemos,sec:eval-audit}). Two results from building them generalize: whole-region
rewrites are necessary only when placement within a region is deterministic
(\cref{lem:slot}), and a partition index costs nothing in leakage once it is logical rather than
physical (\cref{sec:mnemos-m}).
\item A leakage ledger that states every measured channel, for every design and adversary class, as
a lower bound on the information it carries in bits, and every null as the smallest effect the
experiment could have detected, so the design space can be read on one axis
(\cref{sec:eval-ledger}).
\end{enumerate}
